\chapter{What the Calculus Covers}
\label{ch:covers}

This chapter provides a compact, honest summary of what is actually
formalized and proved in the World-Model Calculus, as opposed to merely
discussed or conjectured.  It answers three questions: Which claims are
kernel-checked?  Which desiderata
(Chapter~\ref{ch:desiderata}) are fully discharged?  And where does the
formalization stop?

Every result in this chapter carries a formalization tier:
\begin{description}
\item[\textsc{Tier~A}:] Kernel-checked Lean~4 theorem, zero sorry, zero
  axioms beyond standard foundations.
\item[\textsc{Tier~B}:] Formalized in Lean~4 with sorry-free main
  statement but possible sorry in auxiliary lemmas.
\item[\textsc{Tier~C}:] Proof sketch with full statement in Lean but
  unfinished proof term.
\item[\textsc{Tier~D}:] Open program --- stated precisely, not yet
  formalized.
\end{description}

\noindent
The desiderata D1--D8 serve as the contract the calculus must
fulfill.  The ledger at the end of this chapter
(\S\ref{sec:covers-ledger}) records which desiderata are fully
discharged by Tier~A/B results, which are partially discharged, and
which remain open.


%% ═══════════════════════════════════════════════════════════════
\section{The Kernel}
\label{sec:covers-kernel}
%% ═══════════════════════════════════════════════════════════════

The foundation is a world-model interface with evidence
compositionality: revision of two world-model states and subsequent
query extraction must agree with extracting from each state separately
and then combining the evidence
(Chapter~\ref{ch:world-models}).

\paragraph{World-model interface (\textsc{Tier~A}, \claimCore{}).}
The \code{WorldModel} typeclass defines three operations---\code{revise},
\code{empty}, \code{extract}---over typed parameters
$\mathrm{State}$, $\mathrm{Query}$, $V$.  This discharges
\textbf{D1} (State Primacy): the system carries a revisable state as its
primary representation, not a scalar summary.

\paragraph{Additive regime and profile-hom equivalence
  (\textsc{Tier~A}, \claimCore{}).}
The \code{AdditiveWorldModel} subclass adds the decisive law
$\mathrm{extract}(W_1 + W_2, q) = \mathrm{extract}(W_1, q) +
\mathrm{extract}(W_2, q)$.  The profile-hom equivalence
(Theorem~\ref{thm:universal-property}) shows that additive world models
are exactly answer-profile generators: monoid homomorphisms from states
to answer profiles.  This discharges \textbf{D2} (Compositional
Extraction) in the additive regime.

\paragraph{Dirichlet reference instance (\textsc{Tier~A}, \claimCore{}).}
A concrete Dirichlet instance with query extraction over finite
world-sets instantiates the interface, and the resulting system admits a
sequent-style calculus in which query-revision is an admissible rule.

\paragraph{Local-completeness no-go
  (\textsc{Tier~A}, \claimBoundary{}).}
The no-go theorem (Chapter~\ref{ch:no-go}) demonstrates that no
link-local evidence representation can recover full joint evidence---a
sharp boundary on the ambitions of any PLN-like system.  This is not a
defect of the formalization; it is a theorem about the problem.  This
contributes to discharging \textbf{D8} (Honest Boundaries).

\maplecourt{%
The kernel is visible in every Maple Court scenario: each apartment
carries a state $(n^+, n^-)$ rather than a bare strength scalar; the
building controller revises states additively as sensor reports arrive;
and the no-go theorem explains why the building-level inference cannot
reconstruct the full joint from per-apartment summaries alone.}


%% ═══════════════════════════════════════════════════════════════
\section{Evidence Carriers}
\label{sec:covers-carriers}
%% ═══════════════════════════════════════════════════════════════

Evidence is not one fixed datatype.  The general pattern is that a
carrier stores sufficient statistics for a conjugate family, and revision
in the additive regime is addition of those sufficient statistics.

\paragraph{Three formalized carriers (\textsc{Tier~A}, \claimCore{}).}
Binary evidence $(n^+, n^-)$, Dirichlet evidence
$(\alpha_1, \ldots, \alpha_k)$, and Normal-Gamma evidence
$(n, \sum x_i, \sum x_i^2)$ are each proved to satisfy the
\code{AddCommMonoid} laws under componentwise addition
(Chapter~\ref{ch:evidence}).  All five algebraic rules hold for every
carrier.  This discharges \textbf{D3} (Generic Carriers): the value
type~$V$ is a parameter, not a fixed datatype.

\paragraph{Quantale and Heyting structure
  (\textsc{Tier~A}, \claimConditional{}).}
For the binary carrier specifically, the formalization proves that
evidence forms both a quantale (with a tensor product modelling revision)
and a Heyting algebra (with implication modelling conditional evidence).
These are \emph{binary-specific} structures: the quantale tensor and the
Heyting implication do not generalize to all carriers.  They are
formalized under the side condition that the carrier is binary.


%% ═══════════════════════════════════════════════════════════════
\section{Views and Projections}
\label{sec:covers-views}
%% ═══════════════════════════════════════════════════════════════

The truth-value views---simple truth values, indefinite truth values,
and distributional views---are all formalized as typed projections from
evidence (Chapter~\ref{ch:views}).

\paragraph{Views are not additive (\textsc{Tier~A}, \claimCore{}).}
Strength, confidence, and interval bounds are proved to be non-additive
projections from evidence carriers
(Proposition~\ref{prop:overview-views-not-semantic-core}).  The strength
non-monotonicity theorem shows that more evidence can yield lower
strength.  This discharges \textbf{D4} (Evidence over Views): views are
derived projections, not primary semantic objects.

\paragraph{Walley--Bayesian interval separation
  (\textsc{Tier~B}, \claimConditional{}).}
An explicit separation of Walley predictive intervals from Bayesian
credible intervals is formalized, preventing conflation of two
structurally different interval concepts.

\paragraph{Convergence bridges (\textsc{Tier~B/C}).}
Two convergence bridges ground the asymptotic story:
\begin{itemize}
\item De~Finetti's exchangeability theorem justifies count sufficiency
  (\textsc{Tier~B}).
\item A Solomonoff convergence bridge shows that evidence counts
  approximate the universal posterior (\textsc{Tier~C} --- the main
  statement is in Lean; the convergence argument uses a sorry'd
  measure-theory lemma).
\end{itemize}
A Knuth--Skilling valuation-algebra bridge connects the evidence algebra
to information-theoretic valuations (\textsc{Tier~B}).


%% ═══════════════════════════════════════════════════════════════
\section{State Operations}
\label{sec:covers-state-ops}
%% ═══════════════════════════════════════════════════════════════

State operations---forgetting, provenance tracking, and replay---belong
near the kernel because they demonstrate what a posterior state can do
that a bare truth value cannot: it can be updated, partially retracted,
audited, and explained (Chapter~\ref{ch:state-ops}).

\paragraph{Forgetting layer (\textsc{Tier~A}, \claimCore{}).}
The \code{ForgettingLayer} structure provides scoped retraction: for
every scope~$S$ and state~$W$, forgetting preserves all answers
outside~$S$.  This discharges \textbf{D6} (Retractability).

\paragraph{Conservation theorems (\textsc{Tier~A}, \claimCore{}).}
The \code{EvidenceConservationPack} proves three conservation laws under
exact-inverse forgetting: anti-hallucination, Noether conservation, and
zero leakage.  The exact-inverse no-go theorem (\claimBoundary{}) shows
that if evidence leaked, forgetting cannot exactly undo it.  Together
these discharge \textbf{D7} (Evidence Conservation).

\paragraph{Provenance tracking (\textsc{Tier~A}, \claimCore{}).}
Provenance-tagged revision records the source of each evidence
contribution, enabling the forgetting layer to retract specific sources
while preserving others.

\maplecourt{%
Three hallway humidity alerts arrive, then sensor~3B-north is found to
be miscalibrated.  The forgetting layer retracts $(3, 1)$ from
sensor~3B-north, preserving the $(2, 0)$ from sensor~3B-south.  The
conservation theorem guarantees that the building-level humidity answer
for apartments other than 3B is unchanged.  Without provenance (D6) and
conservation (D7), this surgical retraction is impossible.}


%% ═══════════════════════════════════════════════════════════════
\section{Compiled Inference}
\label{sec:covers-compiled}
%% ═══════════════════════════════════════════════════════════════

The extensional inference rules of classical PLN---deduction,
induction, abduction---are proved exact for the chain, fork, and
collider BN motifs under explicit screening-off and positivity
conditions (Chapter~\ref{ch:bn-models}).

\paragraph{BN exactness theorems (\textsc{Tier~A}, \claimConditional{}).}
D-separation discharge is packaged as a typeclass, so that a user need
only declare the BN structure and the side conditions are checked
automatically.  The compiled deduction formula agrees with full posterior
extraction when and only when the screening-off side condition holds.

\paragraph{$\Sigma$-guarded rewrite system
  (\textsc{Tier~A}, \claimConditional{}).}
The rewrite system (Chapter~\ref{ch:sigma-rewrites}) collects these
compiled rewrites with their typed BN constructor lifts, and a one-call
pipeline composes selector, rewrite, and threshold into a single
certified step.  Every compiled rule carries its structural assumptions as
a typed, machine-checkable predicate~$\Sigma$.  This discharges
\textbf{D5} (Explicit Side Conditions): when $\Sigma$ holds, the
shortcut is exact; when $\Sigma$ fails, the shortcut is not licensed.

\paragraph{Threshold transport (\textsc{Tier~A}, \claimConditional{}).}
Threshold transport with explicit error boundaries
(Chapter~\ref{ch:error-transport}) ensures that confidence gates
compose correctly across inference chains.

\paragraph{Five chaining-limitation no-go theorems
  (\textsc{Tier~A}, \claimBoundary{}).}
Coverage collapse, Lipschitz sensitivity amplification, variance
accumulation, topology--CPT gap, and regime underdetermination
characterize the fundamental limitations of composing per-step
certificates.  These are first-class contributions to \textbf{D8}
(Honest Boundaries).

\maplecourt{%
The compiled chain rule for
$\mathrm{PipeLeak} \to \mathrm{WallHumidity} \to \mathrm{MoldRisk}$
carries side condition $\Sigma = \{\mathrm{ScreenOff}, 0 < s_B < 1\}$.
The $\Sigma$-discipline (D5) ensures the rule fires only when the BN
structure guarantees d-separation.  The chaining-limitation theorems
(D8) tell the building controller that propagating only scalar
summaries through longer chains inevitably loses information.}


%% ═══════════════════════════════════════════════════════════════
\section{Semantic Extensions}
\label{sec:covers-extensions}
%% ═══════════════════════════════════════════════════════════════

Three semantic families extend the kernel beyond extensional
propositional inference:

\begin{itemize}
\item \textbf{Intensional inheritance}
  (Chapter~\ref{ch:intensional}, \textsc{Tier~A/B},
  \claimConditional{}) is grounded through typed
  world-model channels---extensional, ASSOC, and PAT---with the
  Solomonoff log-ratio bridge as the semantic anchor.  A constructive
  no-go (\claimBoundary{}) rules out collapsing the three-channel mixed
  rule back to the original binary law.

\item \textbf{Temporal inference}
  (Chapter~\ref{ch:temporal}, \textsc{Tier~A/B},
  \claimConditional{}) is formalized as another world-model instance
  rather than a separate ontology.  Temporal evidence revision and the
  temporal quantale with free transitivity theorems are Tier~A.  The
  causal-intervention connection to do-calculus is stated as a
  typed interface (\textsc{Tier~C}) --- the full do-calculus
  formalization remains an open program.

\item \textbf{Normative and governance reasoning}
  (Chapter~\ref{ch:governance}, \textsc{Tier~B/C},
  \claimConditional{}) demonstrates that deontic modalities,
  preference aggregation, and governance constraints can be expressed
  within the world-model framework.  The graded deontic evidence
  structure is Tier~B; the temporal-deontic bridge is Tier~C.
\end{itemize}

\paragraph{Higher-order reduction (\textsc{Tier~B/C}).}
Higher-order reduction is formalized with corrected side conditions.
The typed formulation (Chapter~\ref{ch:typed}) provides
root-sum-of-squares bounds under explicit assumptions (\textsc{Tier~B}).
Finite-model quantifier and fuzzy quantifier semantics are proved in the
weakness-based formulation (\textsc{Tier~B}).  The full 13-axis vertex
family and OSLF modal operators are partially formalized
(\textsc{Tier~C}).


%% ═══════════════════════════════════════════════════════════════
\section{Large-Scale Inference and Control}
\label{sec:covers-scale}
%% ═══════════════════════════════════════════════════════════════

The formalization addresses how the calculus scales beyond single-step
inference (Chapter~\ref{ch:scale-control}).

\paragraph{Multideduction (\textsc{Tier~A}, \claimConditional{}).}
Multideduction is formalized via a residual decomposition that makes
the inclusion-exclusion error explicit, together with an
identifiability no-go (\claimBoundary{}) showing that the residual
cannot in general be recovered from local information.

\paragraph{Trail-free protocol (\textsc{Tier~A}, \claimBoundary{}).}
The trail-free deterministic update protocol admits a counterexample
(a 2-cycle that never stabilizes), but damped convergence under
eventual fresh evidence is proved.  Pooling uniqueness under external
Bayesianity pins down the fusion rule.

\paragraph{Inference control (\textsc{Tier~B}, \claimConditional{}).}
BRGI optimization with feasible-set characterization provides the
scheduling foundation.  Premise-selection coverage is proved
submodular, yielding the classical greedy $(1 - 1/e)$ approximation
guarantee, while a coverage surrogate no-go (\claimBoundary{}) shows
that no purely local surrogate can faithfully approximate global
coverage.


%% ═══════════════════════════════════════════════════════════════
\section{Cross-Cutting Bridges}
\label{sec:covers-bridges}
%% ═══════════════════════════════════════════════════════════════

The formalization establishes bridges to eight neighboring inference
frameworks (Chapter~\ref{ch:bridges}):
\begin{enumerate}
\item Markov Logic Networks --- log-linear embedding with profile-hom
  correspondence (\textsc{Tier~A}, \claimCore{}).
\item ProbLog --- distribution-semantics embedding under acyclicity
  (\textsc{Tier~A}, \claimConditional{}).
\item Probabilistic programming languages (Church, Stan, Pyro) ---
  sufficient-statistic embedding under exchangeability
  (\textsc{Tier~A}, \claimConditional{}).
\item PLN v0.9 --- all seven classical rules as $\Sigma$-guarded
  compiled rewrites (\textsc{Tier~A}, \claimConditional{}).
\item NARS --- frequency-confidence bijection with binary evidence
  (\textsc{Tier~A}, \claimCore{}).
\item Classical FOL / set theory --- Boolean degenerate instance
  (\textsc{Tier~A}, \claimCore{}).
\item Markov categories --- profile-hom as commutative-monoid algebra
  in a Markov category (\textsc{Tier~A}, \claimCore{}).
\item $\rho$-calculus --- parallel composition as revision, barbs as
  extraction (\textsc{Tier~A/B}, \claimCore{}).
\end{enumerate}

\noindent
These are exact bridges, not analogies: each neighboring framework is
shown to be a special case of the world-model calculus under explicit
structural assumptions, and the assumptions are stated precisely
enough that violations can be detected and diagnosed.  MeTTa parity
theorems verify executable formula parity between Lean theorems and
MeTTa truth functions.


%% ═══════════════════════════════════════════════════════════════
\section{Desiderata Discharge Ledger}
\label{sec:covers-ledger}
%% ═══════════════════════════════════════════════════════════════

The eight desiderata of Chapter~\ref{ch:desiderata} constitute the
contract the calculus must fulfill.  The following ledger records the
discharge status of each.

\begin{center}
\begin{longtable}{clp{48mm}cl}
\toprule
\textbf{ID} & \textbf{Name} & \textbf{Discharge site}
  & \textbf{Tier} & \textbf{Status} \\
\midrule
\endfirsthead
\toprule
\textbf{ID} & \textbf{Name} & \textbf{Discharge site}
  & \textbf{Tier} & \textbf{Status} \\
\midrule
\endhead
D1 & State Primacy
  & \code{WorldModel} typeclass (Ch~\ref{ch:world-models})
  & A & \textsc{Full} \\
D2 & Composition
  & Profile-hom equivalence (Ch~\ref{ch:world-models})
  & A & \textsc{Full} \\
D3 & Generic Carriers
  & Three carriers (Ch~\ref{ch:evidence})
  & A & \textsc{Full} \\
D4 & Evidence $>$ Views
  & Lossy projections (Ch~\ref{ch:views})
  & A & \textsc{Full} \\
D5 & Explicit $\Sigma$
  & $\Sigma$-guarded rewrites (Ch~\ref{ch:sigma-rewrites})
  & A & \textsc{Full} \\
D6 & Retractability
  & Forgetting layer (Ch~\ref{ch:state-ops})
  & A & \textsc{Full} \\
D7 & Conservation
  & Conservation pack (Ch~\ref{ch:state-ops})
  & A & \textsc{Full} \\
D8 & Honest Boundaries
  & No-go catalog (Chs~\ref{ch:no-go},
    \ref{ch:error-transport}, \ref{ch:scale-control})
  & A & \textsc{Full} \\
\bottomrule
\end{longtable}
\end{center}

\noindent
All eight desiderata are discharged at Tier~A (kernel-checked) for
the additive regime.  Three caveats qualify this claim:
\begin{enumerate}
\item \textbf{D2} is discharged only in the additive regime.  The
  overlap-aware, trust-gated, and coalition regimes do not satisfy
  additive compositionality and are not covered by the current
  discharge (Chapter~\ref{ch:world-models}).
\item \textbf{D3} is discharged for three carriers (binary, Dirichlet,
  Normal-Gamma).  The amplitude carrier $V = \mathbb{C}$ is formalized
  as an additive monoid but does not yet have a full reference instance
  (\textsc{Tier~C}).
\item \textbf{D8} is discharged by a substantial catalog of no-go
  theorems, but completeness of the boundary catalog is not itself
  proved---there may be additional limitations not yet identified.
\end{enumerate}


%% ═══════════════════════════════════════════════════════════════
\section{What the Calculus Does Not Cover}
\label{sec:covers-stops}
%% ═══════════════════════════════════════════════════════════════

Honest boundaries (\textbf{D8}) require stating not only impossibility
theorems but also the gaps in the formalization itself.

\begin{description}
\item[Non-additive regimes (\textsc{Tier~D}).]
  The overlap-aware, trust-gated, and coalition revision regimes are
  typed but not developed in full generality.  No profile-hom
  equivalence is known outside the additive regime.

\item[Measure-theoretic grounding (\textsc{Tier~D}).]
  The Giry monad connection is stated but the terminal
  $\code{WorldModel (Measure\;\Omega)\;Query}$ instance remains
  future work.  The current formalization works with finitary carriers
  and discrete state spaces.

\item[Higher-order completeness (\textsc{Tier~D}).]
  The calculus handles quantified and higher-order queries as instances
  of the generic Query type, but no completeness theorem is proved for
  the full higher-order fragment.

\item[Computational complexity (\textsc{Tier~D}).]
  The calculus says \emph{what} the correct answer is but does not
  provide complexity bounds for computing it.  Exact BN inference is
  \#P-hard in general; the WM calculus licenses compiled shortcuts but
  does not analyze their computational cost.

\item[The probability hypercube (\textsc{Tier~C}).]
  The full classification of theories by which combination of merge-law
  and carrier properties they satisfy is a work in progress.  The
  29-file Lean formalization maps the axes but does not yet prove the
  complete classification theorem.

\item[Full do-calculus formalization (\textsc{Tier~C}).]
  The temporal world-model instance and temporal quantale are
  Tier~A.  The causal-intervention connection to do-calculus is
  stated as a typed interface but not fully formalized.

\item[GSLT universality (\textsc{Tier~C}).]
  The multiset-ensemble canonical world model is Tier~A.  The full GSLT
  universality conjecture---that every graph-structured lambda theory
  induces an additive world model---is stated but not yet proved.
\end{description}


%% ═══════════════════════════════════════════════════════════════
\section{Artifact and Reproducibility}
\label{sec:artifact}
%% ═══════════════════════════════════════════════════════════════

The complete formalization is implemented in Lean~4 with Mathlib
dependencies.  All Tier~A theorems are kernel-checked with zero axioms
beyond the standard foundations.

\paragraph{Reproducing the Lean proofs.}

\begin{lstlisting}
cd lean-projects/mettapedia
lake build
\end{lstlisting}

\noindent
A sorry audit (\texttt{grep -rn sorry} over the relevant modules)
should return zero for all Tier~A modules.

\paragraph{Audit checklist.}
\begin{itemize}
\item \textbf{Lean toolchain:} version and Mathlib pinned in
  \texttt{lean-toolchain} and \texttt{lakefile.lean}.
\item \textbf{Sorry audit:} all Tier~A modules sorry-free.
\item \textbf{Tier labels:} every theorem in the book carries a tier
  label (A/B/C/D) so the reader can verify the formalization status
  independently.
\item \textbf{Holdout discipline:} every benchmark uses fixed
  train/test splits with no hyperparameter tuning on test sets.
\item \textbf{Tagged release:} a tagged commit with CI logs
  accompanies the final version.
\end{itemize}

\maplecourt{%
A reviewer auditing the Maple Court claims can verify each one:
\texttt{lake build} checks the kernel theorems (D1--D2, D5--D7
discharge); \texttt{grep -rn sorry} flags any incomplete proofs;
the tier labels in this chapter map each claim to its formalization
status.  The desiderata ledger (\S\ref{sec:covers-ledger}) shows
that all eight desiderata are discharged at Tier~A for the additive
regime, with three explicit caveats for the frontier directions.}
